% ============================================================
% config/04_ornements.tex
% Motifs ornementaux (pgfornament) - POINT UNIQUE DE MODIFICATION
% Identifiants issus du catalogue du package pgfornament (voir
% "texdoc pgfornament" pour la planche complete des motifs disponibles).
% Changer un identifiant ici suffit a changer le motif partout ou il
% est utilise (page de garde, en-tetes, pieds de page).
% ============================================================

% --- Motifs de la bordure de page de garde ---
\newcommand{\agtOrnCorner}{24}  % motif de coin
\newcommand{\agtOrnSide}{19}    % motif de bordure repete (haut/bas)
\newcommand{\agtOrnSize}{10mm}  % taille d'un motif de bordure

% --- Motifs legers d'en-tete / pied de page (chaque page du corps) ---
\newcommand{\agtOrnHeader}{80}
\newcommand{\agtOrnFooter}{83}
\newcommand{\agtOrnHeaderSize}{5mm}
\newcommand{\agtOrnFooterSize}{4.5mm}

% --- Motif de la page de dedicace (reutilise le motif de coin,
% \agtOrnCorner, deja valide a la compilation sur la page de garde,
% plutot qu'un nouvel identifiant non teste) ---
\newcommand{\agtOrnDedicaceSize}{16mm}

% --- Bordure complete (page de garde uniquement) ---
% CORRECTION : l'ancienne version (boucle \foreach avec rotation
% manuelle des noeuds) faussait le calcul des ancres et produisait des
% blocs pleins noirs au lieu du motif repete. Remplacee par la technique
% eprouvee du template ENSPY (Yuhala, depot latex-thesis) : une chaine
% de noeuds TikZ (bibliotheque "chains") qui tourne aux quatre coins,
% bien plus fiable que des calculs de position manuels.
% LIMITE ANTICIPEE : le nombre de motifs par cote (17 en haut/bas, 25
% sur les cotes) est calibre pour une page A4 avec un motif de 10mm ;
% si la taille du motif (\agtOrnSize) change, ces nombres doivent etre
% ajustes en consequence pour ne pas deborder ou laisser un espace vide.
\newcommand{\pageborderornament}{%
  \begin{tikzpicture}[remember picture, overlay, start chain, node distance=-2mm, color=reportprimary]%
    \node (agtnw) [shift={(5mm,-5mm)}, anchor=north west, on chain] at (current page.north west) {\pgfornament[width=\agtOrnSize]{\agtOrnCorner}};%
    \foreach \i in {1,...,16} {\node [on chain] {\pgfornament[width=\agtOrnSize]{\agtOrnSide}};}%
    \node (agtne) [on chain] {\pgfornament[width=\agtOrnSize]{\agtOrnCorner}};%
    \foreach \i in {1,...,24} {\node [continue chain=going below, on chain] {\pgfornament[width=\agtOrnSize]{\agtOrnSide}};}%
    \node (agtse) [on chain] {\pgfornament[width=\agtOrnSize]{\agtOrnCorner}};%
    \foreach \i in {1,...,16} {\node [continue chain=going left, on chain] {\pgfornament[width=\agtOrnSize]{\agtOrnSide}};}%
    \node (agtsw) [on chain] {\pgfornament[width=\agtOrnSize]{\agtOrnCorner}};%
    \foreach \i in {1,...,24} {\node [continue chain=going above, on chain] {\pgfornament[width=\agtOrnSize]{\agtOrnSide}};}%
  \end{tikzpicture}%
}

% --- Lettrine coloree (raccourci) ---
% Usage dans un chapitre : \agtlettrine{P}{our mener a bien...}
\setcounter{DefaultLines}{3}
\newcommand{\agtlettrine}[2]{\lettrine[lines=3]{\color{reportprimary}#1}{#2}}